\subsection{Herangehensweise}

Der erste Teil der Arbeit behandelt die Zielstellung sowie die Abgrenzung der Problemstellung. Auch werden technische Grundlagen, die für das Verständnis der späteren Kapitel notwendig sind, erklärt. Dazu zählt ein Überblick über die Arbeitsweise eines digitalen Prozessors, sowie eine Erklärung des Konzeptes ``ISA''.

Anschließend folgt eine Erklärung des Konzepts und wie dieses entwickelt wurde. Insbesondere auf den Entwurf der ISA sowie des Befehlssatzes wird detailiert eingegangen. Auch die zugrunde liegenden Designprinzipien, auf welchen die Entscheidungsfindung beruht, werden erläutert. Es folgt ein Überblick über die Struktur der CPU. Dieser Überblick erfolgt nach dem Bottom-Up-Prinzip, angefangen bei individuellen logischen Gattern, über daraus gebildete funktionale Module, bis hin zur Verbindung dieser Module zu einer voll funktionsfähigen CPU.

Nach der Vorstellung des Konzepts gibt es eine Dokumentation der Implementierung. Diese befasst sich mit Implementierungsdetails, welche insbesondere die elektrotechnischen Feinheiten beleuchten. Zudem wird auf durchgeführte Simulationen eingegangen, welche verwendet werden, um den logischen Entwurf alleinstehend zu verifizieren und um die zugehörigen Software-Tools zu testen und entwickeln. Auch die Entwicklung dieser Tools wird beschrieben; darunter fällt hauptsächlich die Entwicklung eines Assemblers für die ISA, sowie das Verfassen von Testprogrammen, welche auf der CPU ausgeführt werden können.

Nach der Dokumentation der Implementierung folgt ein Kapitel über die Inbetriebnahme. Dieses behandelt vor Allem den physikalischen Aufbau der CPU-Module und deren Zusammenschaltung. Zudem wird das Testen und Debuggen der Module und der fertigen CPU dokumentiert.

Ein rückblickendes Fazit bildet den Abschluss der Arbeit. Dabei werden die Testergebnisse, Leistungsanalysen und Funktionstests bewertet. Schwierigkeiten und Probleme, die während des Projekts auftreten, werden ebenfalls aufgeführt. Es folgt ein Ausblick, um mögliche Verbesserungen oder Erweiterungen zu erfassen und Absprungpunkte für künftige Arbeiten abzugrenzen. Abschließend werden persönliche Erkenntnisse dokumentiert.
